% FANET-TopoGNN independent R7 methodological addendum -- NOT SUBMITTED
% This text is proposed for a future, methodologically compatible rerun.
% Do not merge pilot numbers into historical confirmatory tables.
\subsection{Forecast eligibility and full-timeline event accountability}
For a forecast horizon $h$, a decision at time $t$ is labelled only if the
corresponding true observation at $t+h$ exists in the same complete run. The
final $h$ observations remain in the full physical trajectory for event-count
and exposure accounting, but do not acquire duplicated terminal labels. We
separate the forecast-eligible decision view from the full observed trajectory
and map predictions by immutable run and time identifiers. In particular,
missing forecasts are not silently shifted or imputed. The full-timeline event
matching protocol counts connected-to-fragmented transitions over the entire
observed test trajectory, while alarms are emitted only at actually scored
and eligible times. Consequently, previously reported event scores obtained
from truncated timelines and scores from this full-timeline protocol must not
be presented as directly comparable confirmatory measurements.

\subsection{Validation roles and model-family comparison}
Validation groups are separated by the common source trajectory, not merely by
radio scenario or graph-policy variant. Neural training uses disjoint
early-stopping, probability-calibration, and alarm-threshold selection groups,
with final test trajectories untouched. Static and temporal predictors must
be evaluated over an identical set of $(\mathrm{run\_id},t)$ keys and labels;
models with missing test outputs cannot be silently compared. A comparison
across neural and classical families additionally requires matched validation
freedom and operating-point selection. Common test support alone does not
establish such fairness; classical methods requiring a new matching protocol
must be rerun before formal family-wide claims are made.

\subsection{Auditability and experimental scope}
Complete neural evidence requires model architecture settings, PyTorch weights,
validation role assignments, calibration parameters, selected thresholds,
per-frame predicted component counts, final risk probabilities, true labels,
and cryptographic file checksums. Independent replay must reconstruct
prediction outputs from the saved weights and original trajectories rather
than verifying only aggregate event counts. A one-seed, one-epoch engineering
smoke test demonstrating correct replay is not a replacement for a frozen
20-seed, 100-epoch experiment and cannot be used to estimate between-seed
uncertainty or comparative significance.

\subsection{Controller feasibility and limitations}
Historical packet-level policy evaluation used frozen mobility and precomputed
risk trajectories. Reported packet delivery changes do not constitute
proof of closed-loop radio--mobility operation or field deployment. Reporting
must retain the originally declared 25\% intervention duty and
300\,m\,min$^{-1}$ relay-movement budgets. Policies violating either budget
must not be described as budget-feasible; non-significant Holm-adjusted
comparisons must not be called established superiority. A complete release
requires reproducible packet-event simulation, verified source-version lineage,
and independent paired-seed inference under the same full-timeline scoring
protocol.
